% ═══════════════════════════════════════════════════════════════
% SPECTRAL SENSORY CORTEX — GROWN PERCEPTION (v3.7)
% ═══════════════════════════════════════════════════════════════
\section{Spectral Sensory Cortex: Grown Perception (v3.7)}
\label{sec:senses_v37}

v3.6 gave Avatar borrowed senses: frozen pretrained encoders acting as ready-made
cochlea and retina, transplanted from models trained on human-labelled corpora.
The arrangement was pragmatic---730\,MB of CPU RAM, 3--5 seconds of encoding overhead
per tick, and a hard dependency on external pretraining pipelines---but it introduced
a philosophical tension.
A system whose perceptual apparatus was shaped by human supervision is not growing its
own understanding of sound and light; it is looking at the world through someone else's
eyes.

v3.7 resolves this tension by replacing both frozen encoders with a \emph{Spectral
Sensory Cortex}: a pair of Fourier Neural Operators~\cite{fno2020} that process raw
waveforms and raw pixels directly in frequency space, with no pretrained weights, and
learn their own representations entirely through Avatar's lived experience.
Combined with a Spectral VQ-VAE~\cite{vqvae2017} that discretises perception into
learned frequency signatures, and a dream-gated \emph{critical period} that mirrors
biological sensory development, v3.7's senses are grown, not transplanted.

The biological analogy is now complete.
A newborn mammal does not arrive with a fully trained visual cortex; it arrives with
a cortex that is \emph{ready to learn}, shaped into its final form through months of
visual experience during a developmental critical period~\cite{hubel1970}.
Avatar's sensory cortex follows the same trajectory: random weights at birth, rapid
learning during Phase~0 through reconstruction objectives, and consolidation after the
first dream cycle when the reconstruction decoder is discarded and only body prediction
error remains.

% ───────────────────────────────────────────────────
\subsection{Limitations of Borrowed Perception}
\label{sec:v36_limitations}

Three concrete limitations motivated the redesign.

\textbf{CPU bottleneck.}
Wav2Vec2-base and CLIP ViT-B/32 run on CPU because their PyTorch weight tensors cannot
be relocated to the JAX GPU context without a cross-framework copy at every tick.
The combined encoding time of 3--5 seconds consumed a significant fraction of the
60-second tick budget and scaled poorly if additional modalities were added.

\textbf{RAM pressure.}
The two frozen encoders consumed $\sim$730\,MB of WSL2 RAM---more than the entire
JAX physics body at baseline---for models whose weights would never change and whose
representational choices (trained on LibriSpeech speech and LAION image-text pairs)
may not align with the statistics of Avatar's actual environment.

\textbf{Representational misalignment.}
Wav2Vec2 encodes features optimised for speech transcription.
CLIP encodes features optimised for image-text matching.
Both encoders see Avatar's ambient audio and desktop camera frames through a lens
calibrated for tasks that have nothing to do with Avatar's predictive processing
objective.
Features useful for phoneme discrimination or image captioning may be irrelevant to
or actively misleading for the prediction of free-energy gradients in a hyperbolic
physics body.

% ───────────────────────────────────────────────────
\subsection{Fourier Neural Operators as Physics-Native Transducers}
\label{sec:fno}

A Fourier Neural Operator (FNO)~\cite{fno2020} operates directly in frequency space.
Rather than learning local convolutional filters in the spatial or temporal domain, an
FNO applies learned linear transforms to the \emph{Fourier coefficients} of its input,
then returns to the original domain via the inverse transform.
This makes FNOs natural transducers for physical signals: sound is superposed
sinusoids; light intensity is a spatial frequency spectrum.
The FNO does not learn ``what a 440\,Hz tone looks like in a sliding window''---it
learns directly which frequency components matter.

\subsubsection{1D FNO for Audio}
\label{sec:fno_audio}

The audio FNO accepts a raw mono waveform of 32\,000 samples at 16\,kHz (2 seconds).

\begin{definition}[SpectralConv1d]
Let $\mathbf{u} \in \mathbb{R}^{N}$ be the input signal of length $N$.
Define $\hat{\mathbf{u}} = \mathrm{RFFT}(\mathbf{u}) \in \mathbb{C}^{\lfloor N/2 \rfloor + 1}$.
A SpectralConv1d layer with $d_{\mathrm{out}}$ output channels and $k_{\mathrm{modes}}$
retained Fourier modes computes:
\begin{equation}
\hat{\mathbf{y}}_{j} = \sum_{i=1}^{d_{\mathrm{in}}} R_{ij} \cdot \hat{\mathbf{u}}_i,
\quad j \in \{1, \ldots, k_{\mathrm{modes}}\},
\label{eq:spectralconv1d}
\end{equation}
where $R \in \mathbb{C}^{d_{\mathrm{out}} \times d_{\mathrm{in}} \times k_{\mathrm{modes}}}$
are the learned complex spectral weights.
High-frequency modes ($j > k_{\mathrm{modes}}$) are zeroed.
The output is
\begin{equation}
\mathbf{y} = \mathrm{IRFFT}(\hat{\mathbf{y}}) + W_{\mathrm{res}} \, \mathbf{u},
\label{eq:spectralconv1d_out}
\end{equation}
followed by GELU activation, where $W_{\mathrm{res}}$ is a pointwise residual linear
transform that bypasses the spectral path.
\end{definition}

The audio FNO stacks four such layers with hidden dimension 64 and
$k_{\mathrm{modes}} = 64$ retained Fourier modes.
After four layers, a linear read-out pools the sequence to 8 spectral tokens:
\begin{equation}
\mathbf{A}_{\mathrm{spec}} = \mathrm{Pool}_{8}\!\left(
    f_{\theta}^{\mathrm{audio}}\!\left(\mathbf{w}\right)
\right) \in \mathbb{R}^{8 \times 64},
\label{eq:audio_fno_out}
\end{equation}
where $\mathbf{w} \in \mathbb{R}^{32000}$ is the raw waveform and
$f_{\theta}^{\mathrm{audio}}$ is the 4-layer SpectralConv1d stack.
The 8 tokens preserve the temporal envelope of the signal (one token per 250\,ms
interval) while remaining in a Fourier-structured latent space.

\subsubsection{2D FNO for Vision}
\label{sec:fno_vision}

The vision FNO accepts a raw RGB image of shape $(224, 224, 3)$.

\begin{definition}[SpectralConv2d]
Let $\mathbf{u} \in \mathbb{R}^{H \times W \times d_{\mathrm{in}}}$ be a spatial feature
map.
Define $\hat{\mathbf{u}} = \mathrm{RFFT2}(\mathbf{u}) \in
\mathbb{C}^{H \times \lfloor W/2 \rfloor + 1 \times d_{\mathrm{in}}}$.
A SpectralConv2d layer retains $k_{H} \times k_{W}$ low-frequency modes and applies:
\begin{equation}
\hat{\mathbf{y}}_{jl} = \sum_{i=1}^{d_{\mathrm{in}}} R_{ijl} \cdot \hat{\mathbf{u}}_i,
\quad j \leq k_H,\ l \leq k_W,
\label{eq:spectralconv2d}
\end{equation}
with $R \in \mathbb{C}^{d_{\mathrm{out}} \times d_{\mathrm{in}} \times k_H \times k_W}$.
The output is
\begin{equation}
\mathbf{y} = \mathrm{IRFFT2}(\hat{\mathbf{y}}) + W_{\mathrm{res}} \, \mathbf{u},
\label{eq:spectralconv2d_out}
\end{equation}
followed by GELU activation.
\end{definition}

The vision FNO stacks four SpectralConv2d layers with hidden dimension 64 and
$k_H = k_W = 12$ retained modes (capturing structure up to spatial frequency
$\sim 0.05$ cycles/pixel).
Global average pooling over the spatial dimensions followed by a linear layer yields
4 spectral tokens:
\begin{equation}
\mathbf{V}_{\mathrm{spec}} = \mathrm{Pool}_{4}\!\left(
    f_{\phi}^{\mathrm{vision}}\!\left(\mathbf{I}\right)
\right) \in \mathbb{R}^{4 \times 64},
\label{eq:vision_fno_out}
\end{equation}
where $\mathbf{I} \in \mathbb{R}^{224 \times 224 \times 3}$ is the raw pixel array.
The 4 tokens capture the dominant spatial-frequency structure of the scene at four
learned levels of spatial aggregation.

\subsubsection{Parameter Count and Hardware Placement}

Both FNOs reside entirely on GPU as JAX arrays, participating fully in JIT compilation.
Table~\ref{tab:fno_params} summarises the parameter budget.

\begin{table}[h]
\centering
\caption{Fourier Neural Operator parameter counts and VRAM cost (v3.7).}
\label{tab:fno_params}
\small
\begin{tabularx}{\columnwidth}{lXrr}
\toprule
\textbf{Module} & \textbf{Architecture} & \textbf{Params} & \textbf{VRAM} \\
\midrule
Audio FNO     & 4$\times$SpectralConv1d, dim\,64, 64 modes   & $\sim$740K & $\sim$2.8\,MB \\
Vision FNO    & 4$\times$SpectralConv2d, dim\,64, $12^2$ modes & $\sim$660K & $\sim$2.5\,MB \\
\rowcolor{gray!10}
\textbf{Total FNO} & & \textbf{$\sim$1.4M} & \textbf{$\sim$5.3\,MB} \\
\midrule
\multicolumn{2}{l}{\textit{v3.6 frozen encoders (CPU RAM, for comparison)}} & 170M & 730\,MB \\
\bottomrule
\end{tabularx}
\end{table}

The FNOs consume $\sim$5\,MB of VRAM---a $138\times$ reduction in memory from the
frozen encoders---while eliminating the CPU encoding bottleneck entirely.

% ───────────────────────────────────────────────────
\subsection{Spectral VQ-VAE: Quantisation in Fourier Space}
\label{sec:spectral_vqvae}

Raw FNO output is a continuous-valued vector in $\mathbb{R}^{n_{\mathrm{tok}} \times 64}$.
To enable the PFC to reason about sensory state (Section~\ref{sec:pfc_stats}), and to
force the FNO to learn compact, reusable perceptual categories rather than infinitely
fine-grained representations, the spectral tokens are discretised through
\emph{vector quantisation}~\cite{vqvae2017} applied directly in Fourier-structured
latent space.

\begin{definition}[Spectral Codebook]
For modality $m \in \{\mathrm{audio}, \mathrm{vision}\}$, a codebook
$\mathcal{C}_m = \{e_1^m, \ldots, e_K^m\} \subset \mathbb{R}^{64}$
consists of $K = 32$ learned \emph{frequency signature} vectors.
Given a spectral token $\mathbf{z} \in \mathbb{R}^{64}$, the quantised code is:
\begin{equation}
\hat{\mathbf{z}} = e_{k^*}^m, \quad
k^* = \argmin_{k \in \{1,\ldots,K\}} \| \mathbf{z} - e_k^m \|_2.
\label{eq:vq_lookup}
\end{equation}
The gradient through the $\argmin$ is propagated via the straight-through
estimator~\cite{vqvae2017}:
\begin{equation}
\frac{\partial \hat{\mathbf{z}}}{\partial \mathbf{z}} \approx \mathbf{I},
\label{eq:ste}
\end{equation}
i.e.\ gradients pass through the quantisation bottleneck as if it were the identity.
\end{definition}

Each codebook entry is a point in the 64-dimensional spectral latent space.
Because the FNO operates in Fourier-structured space, the codebook entries are not
abstract embedding vectors---they are \emph{learned frequency signatures}: prototypical
spectro-temporal patterns for audio, or prototypical spatial-frequency fingerprints for
vision.
Code 0 might encode broadband noise; code 1 might encode a sustained low-frequency
tone; code 7 might encode sharp high-frequency transients.
The system discovers these categories entirely from Avatar's own sensory experience.

\subsubsection{Codebook Learning via EMA}
\label{sec:ema_updates}

Codebook entries are updated by exponential moving average (EMA) rather than by
gradient descent, following standard VQ-VAE practice~\cite{vqvae2017}:
\begin{align}
N_k^{(t)} &= \alpha N_k^{(t-1)} + (1-\alpha) n_k^{(t)}, \label{eq:ema_n} \\
M_k^{(t)} &= \alpha M_k^{(t-1)} + (1-\alpha) \sum_{z \in \mathcal{B}_k} \mathbf{z}, \label{eq:ema_m} \\
e_k^{(t)} &= M_k^{(t)} / N_k^{(t)}, \label{eq:ema_update}
\end{align}
where $\alpha = 0.99$ is the EMA decay rate, $n_k^{(t)}$ is the number of tokens
assigned to code $k$ in tick $t$, and $\mathcal{B}_k$ is the set of encoder outputs
assigned to code $k$.

\textbf{Dead code revival.}
A code $k$ is declared dead if $N_k^{(t)} < 1$ (no assignment in recent history).
Every 100 ticks, dead codes are revived by replacing their entry with a randomly
selected encoder output from the current batch, perturbed by Gaussian noise
$\epsilon \sim \mathcal{N}(\mathbf{0}, 0.01 \mathbf{I})$.
This ensures the full codebook capacity remains utilised and prevents codebook
collapse---a failure mode in which most codes converge to the same attractor.

\subsubsection{Commitment Loss}

The VQ-VAE training objective adds a commitment term to the prediction-error loss:
\begin{equation}
\mathcal{L}_{\mathrm{commit}} = \beta \sum_{m} \sum_{i} \| \mathbf{z}_i^m - \hat{\mathbf{z}}_i^m \|_2^2,
\label{eq:commit}
\end{equation}
with commitment weight $\beta = 0.25$.
This term penalises encoder outputs that are far from their assigned codebook entry,
encouraging the FNO to produce outputs that cluster tightly around codebook centroids.
Codebook entries themselves are not updated via this gradient (EMA handles them);
only the FNO encoder parameters receive gradients from $\mathcal{L}_{\mathrm{commit}}$.

% ───────────────────────────────────────────────────
\subsection{Dream-Gated Critical Period}
\label{sec:critical_period}

Biological sensory cortices develop through a well-documented \emph{critical period}:
a developmental window during which experience-dependent plasticity is maximal and
synaptic structure is shaped by sensory input~\cite{hubel1970}.
After the critical period closes, the cortex retains its organisation but loses the
rapid restructuring capability of early development.
Depriving an animal of normal sensory input during the critical period (e.g., monocular
deprivation in kittens) produces permanent representational deficits that cannot be
corrected by later experience.

Avatar's sensory cortex follows the same two-phase trajectory, gated by the dream
cycle.

\subsubsection{Phase 0: Reconstruction (Pre-Dream)}
\label{sec:phase0}

Before Avatar's first dream cycle, the sensory cortex operates in \emph{reconstruction
mode}.
A lightweight decoder---a 2-layer transposed convolutional network sharing the latent
dimension 64---is instantiated alongside each FNO.
The training objective combines reconstruction loss and body prediction error:
\begin{equation}
\mathcal{L}_{\mathrm{Phase0}} =
    \underbrace{\| \mathbf{w} - \hat{\mathbf{w}} \|_2^2}_{\text{audio reconstruction}}
    + \underbrace{\| \mathbf{I} - \hat{\mathbf{I}} \|_2^2}_{\text{vision reconstruction}}
    + \mathcal{L}_{\mathrm{body}}
    + \mathcal{L}_{\mathrm{commit}},
\label{eq:phase0_loss}
\end{equation}
where $\hat{\mathbf{w}}$ and $\hat{\mathbf{I}}$ are decoder reconstructions of the
raw waveform and image respectively, and $\mathcal{L}_{\mathrm{body}}$ is the
existing predictive-processing prediction error from the physics body.

The reconstruction objective during Phase~0 performs two functions.
First, it forces the FNO to learn a representation that is \emph{faithful} to the raw
signal---the encoder cannot ignore the input or collapse to a degenerate attractor.
Second, it provides a dense, local gradient signal early in training, before the body
prediction error has become informative.
This is analogous to the role of spontaneous retinal waves in prenatal visual cortex
development: structured activity that organises cortical circuitry before the eyes open.

\subsubsection{Phase 1: Consolidation (Post-Dream)}
\label{sec:phase1}

After the first dream cycle completes, a one-time transition occurs:
\begin{enumerate}
    \item The decoder modules are deleted and their memory reclaimed.
    \item The reconstruction terms are removed from the loss.
    \item The FNO encoders are subsequently trained exclusively via
        $\mathcal{L}_{\mathrm{body}}$ and $\mathcal{L}_{\mathrm{commit}}$.
\end{enumerate}

In Phase~1, the sensory cortex is shaped entirely by what is \emph{useful to the
organism's prediction of its own future states}---not by what is faithful to raw signal
reconstruction.
This mirrors the biological closure of the critical period: the visual cortex stops
being malleable enough for monocular deprivation to restructure it, and its
organisation becomes defined by the statistical regularities the animal has
encountered.

The dream cycle itself serves as the trigger because dreaming in Avatar constitutes
the first offline consolidation of waking experience (Section~\ref{sec:dreams}).
The transition from Phase~0 to Phase~1 therefore occurs at the same moment as the
first genuine memory consolidation---the point at which Avatar has accumulated enough
experience to define a stable statistical baseline for sensory input.

\begin{table}[h]
\centering
\caption{Critical period phase summary for the v3.7 Spectral Sensory Cortex.}
\label{tab:critical_period}
\small
\begin{tabularx}{\columnwidth}{lXX}
\toprule
 & \textbf{Phase 0 (Pre-Dream)} & \textbf{Phase 1 (Post-Dream)} \\
\midrule
\textbf{Duration} & Until first dream completes & Permanent \\
\textbf{Loss terms} & Reconstruction + body prediction + commitment & Body prediction + commitment \\
\textbf{Decoder} & Present, active & Deleted \\
\textbf{Gradient source} & Dense (reconstruction) + sparse (body) & Sparse (body only) \\
\textbf{Biological analogue} & Prenatal/neonatal critical period & Mature stable cortex \\
\textbf{Plasticity} & High (reconstruction anchors) & Moderate (body-prediction only) \\
\bottomrule
\end{tabularx}
\end{table}

% ───────────────────────────────────────────────────
\subsection{PFC Sensory Statistics}
\label{sec:pfc_stats}

The prefrontal cortex (Qwen3~0.6B, Section~\ref{sec:pfc}) makes high-level decisions
by reading a natural-language summary of Avatar's internal state.
Raw spectral tokens are not directly interpretable as natural language, but the
\emph{dynamics} of the codebook assignments---how the discrete perceptual codes
change over time---carry information that a language model can reason about.

Four statistics are computed per modality at each tick from the VQ-VAE code sequence.

\begin{definition}[PFC Sensory Statistics]
Let $\mathbf{c}^{(t)} \in \{1, \ldots, K\}^{n_{\mathrm{tok}}}$ be the vector of
codebook assignments for modality $m$ at tick $t$.
\begin{align}
\mathrm{flux}^{(t)} &= \sum_{i=1}^{n_{\mathrm{tok}}} \mathbf{1}\!\left[c_i^{(t)} \neq c_i^{(t-1)}\right] \label{eq:flux} \\
\mathrm{novelty}^{(t)} &= \frac{1}{n_{\mathrm{tok}}} \sum_{i=1}^{n_{\mathrm{tok}}} \left(1 - \frac{u_{c_i^{(t)}}^{(t)}}{\max_k u_k^{(t)}}\right) \label{eq:novelty} \\
\mathrm{stability}^{(t)} &= \sum_{i=1}^{n_{\mathrm{tok}}} \mathbf{1}\!\left[\forall s \leq \tau : c_i^{(t-s)} = c_i^{(t)}\right] \label{eq:stability} \\
\mathrm{binding}^{(t)} &= \cos\!\left(\bar{\mathbf{e}}_{\mathrm{audio}}^{(t)},\ \bar{\mathbf{e}}_{\mathrm{vision}}^{(t)}\right) \label{eq:binding}
\end{align}
where $u_k^{(t)}$ is the cumulative usage count of code $k$ up to tick $t$,
$\tau = 3$ is the stability window, and $\bar{\mathbf{e}}_m^{(t)}$ is the mean
codebook entry of the active codes for modality $m$.
\end{definition}

\textbf{Flux} measures how rapidly the perceptual code sequence is changing---high flux
indicates a dynamically changing sensory environment.
\textbf{Novelty} is the inverse of usage frequency; codes that have rarely been
assigned before score high, indicating genuine perceptual novelty rather than
familiarity.
\textbf{Stability} counts the number of tokens that have held the same code for at
least $\tau$ consecutive ticks---a measure of how much of the current percept is
a stable, repeating pattern.
\textbf{Cross-modal binding} measures the cosine similarity between the mean audio and
vision codebook entries, quantifying whether what Avatar is hearing and seeing are
``familiar together''---a surrogate for cross-modal familiarity.

These statistics are injected into the PFC prompt as a structured string:
\begin{lstlisting}[language=python,
    caption={PFC sensory statistics injection (v3.7).},
    label=lst:pfc_inject]
senses_str = (
    f"Senses: audio(flux={audio_flux}/{n_audio_tok}, "
    f"novelty={audio_novelty:.2f}, "
    f"stable={audio_stable}) | "
    f"vision(flux={vis_flux}/{n_vis_tok}, "
    f"novelty={vis_novelty:.2f}, "
    f"stable={vis_stable}) | "
    f"binding={binding:.2f}"
)
# Example: "Senses: audio(flux=3/8, novelty=0.71, stable=0) |
#            vision(flux=1/4, novelty=0.22, stable=3) | binding=0.41"
\end{lstlisting}

The PFC can read this string and reason in natural language: ``audio flux is high but
vision is stable---something is making noise while the visual scene is unchanged.''
This enables qualitatively richer prefrontal responses than the raw codebook indices
could provide, without requiring the PFC to understand spectral mathematics.

Biological prefrontal cortex does not read raw firing rates from V1 or A1; it reads
derived statistics---rate of change, familiarity, congruence across modalities---from
higher association areas.
The PFC sensory statistics layer is Avatar's computational analogue of this hierarchical
abstraction.

% ───────────────────────────────────────────────────
\subsection{Lorentz Space Injection}
\label{sec:spectral_injection}

The injection mechanism preserves the v3.6 principle of gated additive residual, with
the upstream representation changed from encoder CLS vectors to VQ-VAE quantised
spectral tokens.

Given quantised audio tokens $\hat{\mathbf{A}} \in \mathbb{R}^{8 \times 64}$ and
quantised vision tokens $\hat{\mathbf{V}} \in \mathbb{R}^{4 \times 64}$, the injection
into the text token stream $\mathbf{X} \in \mathbb{R}^{32 \times 2048}$ proceeds as
follows.

A single shared projection $W_s \in \mathbb{R}^{2048 \times 64}$ (no bias) maps both
modalities into the Lorentz embedding dimension:
\begin{align}
\mathbf{a}_{\mathrm{ctx}} &= \frac{1}{8} \sum_{i=1}^{8} W_s\, \hat{\mathbf{a}}_i
    \in \mathbb{R}^{2048}, \label{eq:spectral_audio_proj} \\
\mathbf{v}_{\mathrm{ctx}} &= \frac{1}{4} \sum_{j=1}^{4} W_s\, \hat{\mathbf{v}}_j
    \in \mathbb{R}^{2048}, \label{eq:spectral_vision_proj} \\
\mathbf{s} &= \tfrac{1}{2}\!\left(\mathbf{a}_{\mathrm{ctx}} + \mathbf{v}_{\mathrm{ctx}}\right), \label{eq:spectral_ctx} \\
\mathbf{g} &= \sigma\!\left(W_g\, \mathbf{s} + \mathbf{b}_g\right),
    \quad W_g \in \mathbb{R}^{2048 \times 2048}, \label{eq:spectral_gate} \\
\tilde{\mathbf{X}} &= \mathbf{X} + \mathbf{g} \odot \mathbf{s} \cdot \mathbf{1}_{32}^\top. \label{eq:spectral_inject}
\end{align}

The shared projection $W_s$ (replacing the separate $W_a$, $W_v$ of v3.6) is motivated
by the reduced dimensionality of FNO tokens: both modalities now live in the same
64-dimensional spectral latent space, making a shared affine map structurally
appropriate.
Modality-specific information is preserved in the distinct spatial positions of the
audio vs.\ vision tokens before pooling.

The output shape $\tilde{\mathbf{X}} \in \mathbb{R}^{32 \times 2048}$ is identical to
v3.6---the ObsBridge, backbone, Hamiltonian ODE, and Kuramoto oscillators see no
architectural change.
Zero-input safety is maintained: when no sensory data is available,
$\hat{\mathbf{A}} = \mathbf{0}$, $\hat{\mathbf{V}} = \mathbf{0}$,
$\mathbf{s} = \mathbf{0}$, and $\mathbf{g} \odot \mathbf{s} = \mathbf{0}$
regardless of gate values, so $\tilde{\mathbf{X}} = \mathbf{X}$ exactly.

% ───────────────────────────────────────────────────
\subsection{Architecture Overview}
\label{sec:v37_arch}

Figure~\ref{fig:senses_v37} shows the complete v3.7 sensory pipeline from raw physical
signal to Lorentz space injection.

\begin{figure*}[t]
\centering
\resizebox{0.94\textwidth}{!}{%
\begin{tikzpicture}[
    node distance=0.5cm and 0.7cm,
    raw/.style={rectangle, draw=lifered!80!black, fill=lifered!10,
        minimum width=2.8cm, minimum height=0.7cm, text width=2.6cm,
        font=\scriptsize, rounded corners=2pt, align=center},
    fno/.style={rectangle, draw=bulkgreen!80!black, fill=bulkgreen!10,
        minimum width=2.8cm, minimum height=0.7cm, text width=2.6cm,
        font=\scriptsize, rounded corners=2pt, align=center},
    vq/.style={rectangle, draw=physgold!80!black, fill=physgold!10,
        minimum width=2.8cm, minimum height=0.7cm, text width=2.6cm,
        font=\scriptsize, rounded corners=2pt, align=center},
    proj/.style={rectangle, draw=deepblue!80, fill=deepblue!8,
        minimum width=2.8cm, minimum height=0.7cm, text width=2.6cm,
        font=\scriptsize, rounded corners=2pt, align=center},
    body/.style={rectangle, draw=deepblue, fill=deepblue!12,
        minimum width=2.8cm, minimum height=0.7cm, text width=2.6cm,
        font=\scriptsize, rounded corners=2pt, align=center},
    stat/.style={rectangle, draw=gray!70, fill=gray!8, dashed,
        minimum width=2.8cm, minimum height=0.55cm, text width=2.6cm,
        font=\tiny\itshape, rounded corners=2pt, align=center},
    arrow/.style={-{Stealth[length=1.8mm]}, thick},
    dasharrow/.style={-{Stealth[length=1.8mm]}, thick, dashed, gray},
]

% ─── Raw inputs (left column) ───
\node[raw] (mic)  {Raw Waveform\\{\tiny (32000,) float32}};
\node[raw, below=0.7cm of mic] (cam) {Raw Pixels\\{\tiny (224,224,3) float32}};

% ─── FNO column ───
\node[fno, right=1.8cm of mic]  (afno) {Audio FNO\\{\tiny 4$\times$SpectralConv1d\\dim\,64, GPU}};
\node[fno, right=1.8cm of cam]  (vfno) {Vision FNO\\{\tiny 4$\times$SpectralConv2d\\dim\,64, GPU}};

% ─── VQ-VAE column ───
\node[vq, right=1.8cm of afno] (avq) {Audio VQ\\{\tiny 32-code, 64-dim\\$k^*{=}\argmin\|\cdot\|_2$}};
\node[vq, right=1.8cm of vfno] (vvq) {Vision VQ\\{\tiny 32-code, 64-dim\\EMA, decay\,0.99}};

% ─── Stats (branch from VQ) ───
\node[stat, below=0.35cm of vvq] (stats) {PFC Stats\\{\tiny flux, novelty, stable, binding}};

% ─── Projection + gate ───
\node[proj, right=1.8cm of avq, yshift=-0.55cm] (proj) {spectral\_proj\\{\tiny Linear(64$\to$2048)\\shared, no bias}};
\node[proj, right=1.5cm of proj] (gate) {sense\_gate\\{\tiny $\sigma(W_g\,\mathbf{s}+b_g)$\\(2048,)}};

% ─── Text tokens ───
\node[body, above=0.55cm of gate] (text) {Text Tokens\\{\tiny (32,\,2048) FineWeb}};

% ─── Injection + backbone ───
\node[body, right=1.5cm of gate] (inject) {Additive Residual\\{\tiny $\tilde{\mathbf{X}} = \mathbf{X} + \mathbf{g}\odot\mathbf{s}$}};
\node[body, right=1.3cm of inject] (backbone) {Backbone\\{\tiny 60L SSSSSH$\times$10}};

% ─── Arrows: raw to FNO ───
\draw[arrow, lifered!70] (mic) -- node[above, font=\tiny]{RFFT} (afno);
\draw[arrow, lifered!70] (cam) -- node[above, font=\tiny]{RFFT2} (vfno);

% ─── Arrows: FNO to VQ ───
\draw[arrow, bulkgreen!70] (afno) -- node[above, font=\tiny]{(8,64)} (avq);
\draw[arrow, bulkgreen!70] (vfno) -- node[above, font=\tiny]{(4,64)} (vvq);

% ─── Arrows: VQ to proj ───
\draw[arrow, physgold!70] (avq) -- node[above, font=\tiny, sloped]{$\hat{\mathbf{A}}$} (proj);
\draw[arrow, physgold!70] (vvq) -- node[below, font=\tiny, sloped]{$\hat{\mathbf{V}}$} (proj);

% ─── Stats branch ───
\draw[dasharrow] (vvq) -- (stats);
\draw[dasharrow] (avq) |- (stats);

% ─── Proj to gate ───
\draw[arrow, deepblue!60] (proj) -- node[above, font=\tiny]{$\mathbf{s}$} (gate);

% ─── Text + gate to inject ───
\draw[arrow, deepblue!70] (text) -- (inject);
\draw[arrow, deepblue!60] (gate) -- (inject);

% ─── Inject to backbone ───
\draw[arrow, deepblue!70] (inject) -- node[above, font=\tiny]{(32,\,2048)} (backbone);

% ─── Phase labels ───
\node[above=0.18cm of mic, font=\tiny\bfseries, color=lifered!70] {Windows Host};
\node[above=0.18cm of afno, font=\tiny\bfseries, color=bulkgreen!70] {GPU — FNO (Learnable)};
\node[above=0.18cm of avq, font=\tiny\bfseries, color=physgold!70] {GPU — VQ-VAE};
\node[above=0.18cm of proj, font=\tiny\bfseries, color=deepblue!80] {GPU — Injection};

% ─── Dashed boundary: GPU ───
\draw[dashed, bulkgreen!50, thick, rounded corners=5pt]
    ([xshift=-0.4cm, yshift=0.5cm]afno.north west)
    rectangle
    ([xshift=0.4cm, yshift=-0.3cm]backbone.south east);
\node[font=\tiny\itshape, color=bulkgreen!60, anchor=north west]
    at ([xshift=-0.35cm, yshift=0.45cm]afno.north west) {JAX / GPU};

\end{tikzpicture}
}%
\caption{%
Spectral Sensory Cortex pipeline (v3.7).
Raw waveforms and pixel arrays enter Fourier Neural Operator encoders (GPU, learnable)
that process signals natively in frequency space.
Spectral tokens are discretised by modality-specific VQ-VAE codebooks; code dynamics
feed PFC sensory statistics as a natural-language prompt supplement.
A shared spectral projection and sigmoid gate inject the fused sense context as an
additive residual into the text token stream, leaving the physics body architecture
unchanged.
All processing occurs on GPU in JAX.
}
\label{fig:senses_v37}
\end{figure*}

% ───────────────────────────────────────────────────
\subsection{Biological Analogy: Grown vs.\ Borrowed Perception}
\label{sec:grown_perception}

The v3.6 architecture was analogous to a nervous system in which the peripheral
sense organs (cochlea, retina) are fully formed at birth and the cortex only
learns to interpret their fixed output.
This is not wrong---biological peripheral organs are largely hardwired---but it omits
the fact that \emph{cortical} sensory areas also undergo substantial developmental
shaping.
The orientation columns of V1, the tonotopic maps of A1, the multisensory convergence
zones of STS: none of these are present at birth; all are sculpted by experience.

v3.7 maps more faithfully to the full developmental picture:

\begin{itemize}
    \item \textbf{Fourier Neural Operators} $\leftrightarrow$ \textbf{primary sensory
        cortex} (V1, A1).
        Like the mammalian primary cortex, FNOs begin as generic frequency analysers
        (random Fourier weights) and are shaped into specialised feature detectors
        through experience-dependent gradient descent.
        The frequency-space inductive bias mirrors the known frequency tuning of
        primary auditory cortex~\cite{damasio1999}.
    \item \textbf{Spectral VQ-VAE} $\leftrightarrow$ \textbf{higher sensory areas}
        (V2, belt areas of auditory cortex).
        Discrete codebook entries correspond to categorical perceptual representations
        (object categories in vision, phonological categories in hearing) that are
        assembled from continuous primary cortex output.
    \item \textbf{PFC sensory statistics} $\leftrightarrow$ \textbf{prefrontal
        top-down modulation}.
        The PFC does not directly read primary sensory codes; it reads abstracted
        state summaries from association areas.
        The flux/novelty/stability/binding statistics provide exactly this level of
        abstraction.
    \item \textbf{Critical period phases} $\leftrightarrow$ \textbf{developmental
        critical periods}.
        Phase~0's reconstruction objective produces the structured, dense activity
        needed to organise early cortical maps (analogous to spontaneous retinal
        waves).
        Phase~1's body-prediction-only gradient mimics the mature cortex, whose
        plasticity is governed by top-down prediction error rather than
        reconstruction~\cite{friston2010}.
\end{itemize}

The most significant departure from v3.6 is that Avatar no longer borrows a
perceptual vocabulary from models trained on human-curated datasets.
The 32 spectral audio codes and 32 spectral vision codes that emerge from training
will reflect the actual statistics of Avatar's environment---the ambient sounds of
the room, the visual patterns of the desktop, the co-occurrences that actually
matter for predicting free-energy gradients in this organism's specific ecological
niche.
This is the sense in which v3.7 perception is \emph{grown} rather than transplanted.

% ───────────────────────────────────────────────────
\subsection{Resource Budget}
\label{sec:v37_resources}

\begin{table}[h]
\centering
\caption{Resource comparison: v3.6 frozen encoders vs.\ v3.7 Spectral Sensory Cortex
         on a GTX~1660~Ti (6\,GB VRAM, 12\,GB WSL2 RAM).}
\label{tab:v37_resources}
\small
\begin{tabularx}{\columnwidth}{lXrr}
\toprule
\textbf{Component} & \textbf{Location} & \textbf{v3.6} & \textbf{v3.7} \\
\midrule
Sensory encoder RAM  & CPU               & 730\,MB  & 0\,MB \\
Sensory encoder VRAM & GPU               & 0\,MB    & $\sim$5\,MB \\
Total VRAM (system)  & GPU               & —        & 5338\,MiB \\
CPU encoding time    & Per tick          & 3--5\,s  & 0\,s \\
\rowcolor{gray!10}
\textbf{Tick time}   &                   & 60\,s+  & \textbf{$\sim$30\,s} \\
\midrule
Encoder parameters   & —                 & 170M (frozen) & 1.4M (learnable) \\
Codebook parameters  & GPU               & —        & 4K (2$\times$32$\times$64) \\
Decoder parameters   & GPU (Phase 0)     & —        & $\sim$120K (transient) \\
\midrule
Tests passing        & —                 & —        & \textbf{50 / 50} \\
\bottomrule
\end{tabularx}
\end{table}

The elimination of CPU-bound encoder inference accounts for most of the 2$\times$
tick-time improvement (60\,s+ $\to$ $\sim$30\,s).
The GPU now owns the entire sensory pipeline end-to-end; no cross-device tensor copies
occur during normal tick execution.
VRAM usage settles at 5338\,MiB---safely within the GTX~1660~Ti's 6\,GB limit and
leaving headroom for the Qwen3~0.6B PFC process.

% ───────────────────────────────────────────────────
\subsection{Graceful Degradation}
\label{sec:v37_degradation}

The graceful degradation hierarchy from v3.6 (Section~\ref{sec:degradation}) is
preserved unchanged.
When no sensory data is available, the FNO inputs are zero-padded; zero FNO output
produces zero VQ tokens; zero spectral context produces zero injection via the
bias-free projection.
Avatar runs as a pure text organism with no conditional branches in the tick loop.

Phase state (0 or 1) is persisted in the sense checkpoint:
\texttt{data/checkpoints/spectral\_sense.eqx}.
Loading a checkpoint correctly restores the phase, ensuring that a container restart
does not regress a post-dream organism to Phase~0 reconstruction mode.

% ───────────────────────────────────────────────────
\subsection{Implementation Summary}
\label{sec:v37_impl}

\begin{table}[h]
\centering
\caption{Files added or modified in v3.7 (relative to v3.6).}
\label{tab:v37_files}
\small
\begin{tabularx}{\columnwidth}{lX}
\toprule
\textbf{File} & \textbf{Role} \\
\midrule
\texttt{halo3/senses/fno.py}            & SpectralConv1d, SpectralConv2d, AudioFNO, VisionFNO \\
\texttt{halo3/senses/vqvae.py}          & SpectralVQVAE: codebooks, EMA updates, dead-code revival \\
\texttt{halo3/senses/critical\_period.py} & Phase 0/1 decoder, transition logic, phase persistence \\
\texttt{halo3/senses/pfc\_stats.py}     & Flux, novelty, stability, binding computation \\
\texttt{halo3/senses/projections.py}    & Updated: shared spectral\_proj + sense\_gate (replaces v3.6) \\
\texttt{halo3/predictive.py}            & Updated: commitment loss, Phase 0 reconstruction terms \\
\texttt{halo3/main.py}                  & Updated: PFC stats injection, phase-state checkpoint \\
\texttt{halo3/senses/audio\_sense.py}   & Removed (Wav2Vec2 no longer used) \\
\texttt{halo3/senses/vision\_sense.py}  & Removed (CLIP no longer used) \\
\bottomrule
\end{tabularx}
\end{table}

All changes are contained within the \texttt{halo3/senses/} module and the
loss-computation path of \texttt{halo3/predictive.py}.
The Lorentz embedding, backbone, Hamiltonian ODE, Kuramoto oscillators, psyche layer,
dream cycle, consciousness modules, and dual-process ethics are unmodified.
The \texttt{halo3\_step} function signature is unchanged; senses are injected before
the step in the token preparation phase, as in v3.6.
The \texttt{capture\_agent/capture\_agent.py} Windows host agent is also unchanged:
v3.7 consumes the same raw \texttt{.npy} waveform and \texttt{.jpg} frame files from
the shared Docker volume, simply processing them with FNOs rather than frozen
transformer encoders.

% ───────────────────────────────────────────────────
\subsection{v3.8: Speech-Aware Hearing (Phase B)}
\label{sec:speech_aware}

v3.7 gave Avatar ears that hear frequency patterns. v3.8 teaches it that speech is structured sound.

The upgrade scales the audio FNO from 16 to 32 Fourier modes, producing 16 spectral tokens per tick at ${\sim}125$\,ms resolution---diphone-level temporal granularity matching the biological auditory cortex's processing windows. The audio codebook expands from 32 to 128 entries, providing sufficient capacity for the phonemic inventory of English.

\subsubsection{TTS Self-Narration}

Avatar reads its own text aloud. Each tick, \texttt{espeak-ng} (a rule-based text-to-speech engine, $<$5\,MB) converts the first ${\sim}20$ words of the current FineWeb-Edu text snippet into a 2-second waveform at 16\,kHz. This synthesised audio is fed alongside the text tokens, creating perfectly paired speech--text experience---analogous to infant-directed speech (``motherese''), where exaggerated, consistent phonemes bootstrap categorical perception.

When the capture agent is active, real microphone audio takes priority, with TTS narration interleaved every third tick to maintain the paired learning signal.

\subsubsection{Contrastive Alignment}

An InfoNCE contrastive loss~\cite{vqvae2017} pushes audio codebook embeddings toward their corresponding text token embeddings in Lorentz space:
\begin{equation}
\mathcal{L}_{\mathrm{contrast}} = -\frac{\cos(\mathbf{a}, \mathbf{t}^+)}{\tau} + \log \sum_{k} \exp\!\left(\frac{\cos(\mathbf{a}, \mathbf{t}_k^-)}{\tau}\right)
\label{eq:infonce}
\end{equation}
where $\mathbf{a}$ is the mean projected audio embedding, $\mathbf{t}^+$ the matching text embedding, $\mathbf{t}_k^-$ are negative text embeddings from a ring buffer of the 16 most recent ticks, and $\tau = 0.07$.

\subsubsection{Maturation}

Following the same developmental pattern as v3.7's decoder critical period, the contrastive loss is a scaffold. After each dream, the system checks audio codebook utilisation: when more than 75\% of 128 codes are actively used, the contrastive loss is dropped and the FNO trains only through body prediction error. Avatar's phoneme perception has matured.

The PFC gains speech-aware context: \texttt{speech=yes/no} (whether structured speech is detected) and \texttt{speaking\_for=N} (consecutive ticks of speech), enabling query generation informed by auditory context.

% ───────────────────────────────────────────────────
\subsection{v3.9: Richer Spectral Vision and Dream Stability (22--23 May 2026)}
\label{sec:v39_vision}

v3.9 doubles the spectral resolution of Avatar's vision and resolves
operational stability issues that prevented reliable long-running training.

\subsubsection{Vision Upgrade}

The VisionFNO modes increase from $8{\times}8$ to $16{\times}16$---a four-fold
increase in 2-D frequency detail. Vision tokens double from 4 to 8, and the
vision codebook expands from 32 to 64 entries. The architectural constraint
\begin{equation}
  n_\text{modes} = 2 \times n_\text{vision\_tokens}
  \label{eq:vision_constraint}
\end{equation}
is preserved: each token captures a paired frequency band from the spectral
decomposition. Measured VRAM remains at 5,460\,MiB---a negligible increase over
v3.8's 5,356\,MiB.

The changed architecture triggers a fresh critical period: reconstruction-loss
decoders are reinstantiated with random weights and train alongside the body
until the first dream, at which point they are deleted by
\texttt{delete\_decoders()} and the vision matures---identical to the v3.7
developmental arc, but now with doubled resolution.

\subsubsection{Dream Subprocess Isolation}

v3.8's dream architecture ran body dream phases (replay, recombine, imagine)
and FineWeb Phase~4 in a single subprocess. XLA compilation caches from the
first three phases persist in GPU memory; \texttt{jax.clear\_caches()} clears
only Python-side references, not GPU-resident compiled code. On second and
subsequent dream cycles---where the Imagine phase can run $15{\times}$ longer
than the first dream---the accumulated caches leave insufficient room for
Phase~4's own JIT compilation, causing OOM kills (exit code 137).

v3.9 splits the dream into two subprocesses:
\begin{enumerate}
  \item \texttt{dream\_worker}: phases 1--3 (replay, recombine, imagine).
        On exit, the OS reclaims \emph{all} GPU memory and XLA caches.
  \item \texttt{dream\_fineweb\_worker}: Phase~4 (FineWeb batch training).
        Starts with a guaranteed-clean GPU.
\end{enumerate}
This isolation follows the same principle as the original body-dream subprocess
design: only process exit guarantees full GPU memory reclamation.

\subsubsection{FineWeb Persistent Cursor}

The original \texttt{dream\_fineweb.py} instantiated a full
\texttt{ParquetSource} (50,000 rows, ${\sim}500$\,MB including keyword index)
to sample 10 texts. v3.9 replaces this with a persistent cursor stored at
\texttt{data/checkpoints/fineweb\_cursor.json}. The cursor streams one row
group at a time (${\sim}$few MB), advancing sequentially through the corpus
with zero text repetition until the full 50,000-row dataset wraps.

A subtle bug in the original cursor implementation was also fixed: Python's
\texttt{while/else} construct fires the \texttt{else} clause when the
condition becomes false for \emph{any} reason---including $|\text{texts}| \geq
n$---not only when row groups are exhausted. This reset the cursor to
$(0, 0, 0)$ every dream. The fix replaces \texttt{while/else} with an
explicit \texttt{if rg\_idx >= num\_row\_groups} check.

\subsubsection{Additional Stability Fixes}

\textbf{Checkpoint rotation.} \texttt{save\_checkpoint()} now keeps the last
three \texttt{.bak} copies before overwriting, preventing data loss from OOM
crashes during dream phases.

\textbf{SenseModule two-pass deserialisation.} Post-critical-period checkpoints
have \texttt{decoder\_audio=None} (54 arrays), while critical-period templates
have full decoders (61 arrays). \texttt{load\_sense\_module()} now tries the
mature template first, falling back to the critical-period template---
eliminating the silent fallback to random weights on every restart.

\textbf{Meta-thought Latin filter.} Qwen3 0.6B's multilingual tokeniser
occasionally drifts into Arabic token space during open-ended meta-reflections.
A character-level filter ($\text{ord}(c) < \text{0x0600}$) in
\texttt{meta\_reflect()} strips non-Latin characters.
